% !TEX root = ../algebra_02.tex
\section{Евклидовы пространства}

\begin{Definition}{Скалярное произведение}{}
	\textit{Скалярное произведение} на $V/\mathbb{R}$ — это симметрическая билинейная форма $\mathcal{B}$ на $V$ такая, что
	\[
	\forall\, v \in V\colon \mathcal{B}(v,v) \geq 0, \quad \mathcal{B}(v,v) = 0 \;\Longleftrightarrow\; v = 0.
	\] 
	(Положительная определённость.) 
\end{Definition}

\begin{Definition}{Евклидово пространство}{}
	\textit{Евклидово пространство} — вещественное линейное пространство $V$ с фиксированным скалярным произведением.
	Обозначение: $\mathcal{B}(v, w) =: (v, w)$ или $\langle v, w \rangle$.
\end{Definition}

Норма вектора: $\|v\| := \sqrt{(v,v)}$. Для неё выполнены:
\[
\|\alpha v\| = |\alpha|\cdot\|v\|, \qquad
|(v,w)| \leq \|v\|\cdot\|w\| \quad \text{(Коши--Буняковский)}, \qquad
\|v+w\| \leq \|v\| + \|w\|.
\] 

\begin{Example}{Евклидово пространство на $C[0,1]$}{}
	$V = C[0,1]$ с $(f,g) = \displaystyle\int_0^1 f(x)\,g(x)\,dx$ — евклидово пространство.
\end{Example}
